\begin{tikzpicture}[sb, x=0.8cm,
    res0/.style={circle, fill=sbRes0, minimum size=4.4pt, inner sep=0pt},
    res1/.style={regular polygon, regular polygon sides=3, fill=sbRes1,
      minimum size=6.4pt, inner sep=0pt},
    res2/.style={diamond, fill=sbRes2, minimum size=5.8pt, inner sep=0pt},
    res3/.style={regular polygon, regular polygon sides=5, fill=sbRes3,
      minimum size=6pt, inner sep=0pt},
    res4/.style={star, star points=5, star point ratio=2.2, fill=sbRes4,
      draw=sbRes3, line width=0.2pt, minimum size=7pt, inner sep=0pt},
    term/.style={sb aux pt, minimum size=6pt},
    coll/.style={sb fail pt, minimum size=7pt},
    lab/.style={font=\footnotesize, inner sep=1pt}]
  \pgfplotstableread{\sbdata{residues.dat}}\sbResTab
  \pgfplotstablegetrowsof{\sbResTab}%
  \pgfmathtruncatemacro\sbResLast{\pgfplotsretval-1}%
  \newcommand\sbResLine{%
    \draw[sbInk, line width=0.6pt] (27.6,0) -- (40.4,0);
    \foreach \k in {28,...,40}{%
      \pgfmathtruncatemacro\r{mod(\k,5)}%
      \draw[sbInk] (\k,-0.08) -- (\k,0.08);
      \node[lab, below] at (\k,-0.08) {$\k$};
      \node[res\r] at (\k,-0.62) {};}}
  \def\gx{42.2}\def\gw{0.75}
  \newcommand\sbResGrid{%
    \foreach \c in {0,...,4}{%
      \node[res\c] at (\gx+\c*\gw+\gw/2,1.45) {};
      \foreach \qq/\yy in {6/0.45,7/0.95}{%
        \draw[sbOther!60, line width=0.3pt]
          (\gx+\c*\gw,\yy-0.25) rectangle ++(\gw,0.5);
        \pgfmathtruncatemacro\nn{5*\qq+\c}%
        \node[font=\tiny, text=sbOther, anchor=north east, inner sep=0.8pt]
          at (\gx+\c*\gw+\gw,\yy+0.25) {$\nn$};}}
    \node[lab, anchor=east] at (\gx-0.05,0.45) {$6$};
    \node[lab, anchor=east] at (\gx-0.05,0.95) {$7$};
    \node[lab, anchor=north, align=center] at (\gx+2.5*\gw,0.12)
      {\scriptsize rows: quotient by $5$\\[-2pt]\scriptsize columns: residue};}
  \node[anchor=west, font=\small] at (27.4,2.55)
    {(a) $a=4<M'=5$: basis $1,x,y,x^2$ of $\calL(4P_0)$,
     pole orders $0,2,3,4$};
  \sbResLine
  \sbResGrid
  \foreach \i in {0,...,\sbResLast}{%
    \pgfplotstablegetelem{\i}{total}\of\sbResTab \let\tot\pgfplotsretval
    \pgfplotstablegetelem{\i}{u}\of\sbResTab \let\uu\pgfplotsretval
    \pgfplotstablegetelem{\i}{poleu}\of\sbResTab \let\pu\pgfplotsretval
    \pgfplotstablegetelem{\i}{polev}\of\sbResTab \let\pv\pgfplotsretval
    \pgfmathtruncatemacro\idx{\i+1}%
    \pgfmathsetmacro\ly{ifthenelse(mod(\i,2)==1,1.15,0.55)}%
    \node[term] at (\tot,0.35) {};
    \ifodd\i \draw[sbAux, line width=0.4pt] (\tot,0.45) -- (\tot,\ly);\fi
    \node[lab, above, align=center, fill=white] at (\tot,\ly)
      {$\uu\,v_{\idx}^{5}$\\[-1pt]{\scriptsize$\pu+5\cdot\pv$}};
    \pgfmathtruncatemacro\rr{mod(\tot,5)}%
    \pgfmathtruncatemacro\qt{(\tot-\rr)/5}%
    \node[term] at (\gx+\rr*\gw+0.24,{0.45+0.5*(\qt-6)-0.09}) {};}
  \draw[sbInk, line width=0.6pt] (38,0.35) circle (0.2);
  \node[lab, anchor=south, align=center] at (38,1.45)
    {unique maximum:\\$-v_{P_0}(w)=38$};
  \node[anchor=west, font=\small] at (27.4,-1.45)
    {(b) $a=5\ge M'$: basis $1,x,y,x^2,xy$ of $\calL(5P_0)$, pole orders
     $0,2,3,4,5$, and $5\equiv0\bmod5$};
  \begin{scope}[yshift=-3.35cm]
    \sbResLine
    \sbResGrid
    \node[coll] at (34.92,0.35) {};
    \node[coll] at (35.08,0.35) {};
    \node[lab, anchor=east, text=sbFail] at (34.65,0.35)
      {$1\cdot v_1^{5}$: $0+5\cdot7$};
    \node[lab, anchor=west, text=sbFail] at (35.35,0.35)
      {$xy\,v_5^{5}$: $5+5\cdot6$};
    \node[lab, anchor=south west, align=left] at (27.8,0.75)
      {same pole order $35$: the leading terms can cancel;
       (N) $a<M'$ excludes this};
    \node[coll] at (\gx+0.2,0.86) {};
    \node[coll] at (\gx+0.32,0.86) {};
  \end{scope}
  \node[font=\footnotesize, anchor=north west, align=left] at (27.4,-4.55) {%
    \tikz\node[term]{}; term $u_iv_i^{5}$, at its pole order
    $-v_{P_0}(u_i)-5\,v_{P_0}(v_i)$\qquad
    \tikz\node[coll]{}; two terms with the same pole order\\[3pt]
    residue mod $5$:\
    \tikz\node[res0]{}; $0$\quad \tikz\node[res1]{}; $1$\quad
    \tikz\node[res2]{}; $2$\quad \tikz\node[res3]{}; $3$\quad
    \tikz\node[res4]{}; $4$\qquad
    \sbschematic[the pole orders of the $v_i$ are illustrative]};
\end{tikzpicture}
\caption{Why the section $w=\sum_iu_iv_i^{M'}$ is nonzero
(Proposition~\ref{prop:independence}, Step~2 of
Theorem~\ref{thm:aux-section}), for $q=25$, $g=1$, $M'=5$. Pole orders
$-v_{P_0}$ are drawn on a number line, with their residues modulo $5$ shown
by shape and grey shade; at the right the same terms are placed in a grid by
quotient and residue modulo $5$, as in Figure~\ref{fig:J-degrees}. (a) For
$a=4$, the adapted basis $1,x,y,x^2$ of $\calL(4P_0)$ has pole orders
$0,2,3,4$. With pole orders $7,6,7,6$ of $v_1,\dots,v_4$, the terms
$u_iv_i^5$ (sky-blue squares) have pole orders $35,32,38,34$, one in each of
four residue classes; the largest is attained once, so it is the pole order
of $w$, and $w\ne0$. A section $w$ computed for one curve with these
parameters has pole order $38$. (b) If $a\ge M'$, here $a=5$, two basis
elements have pole orders congruent modulo $M'$, so two terms can have the
same pole order (purple open circles, one grid cell) and their leading terms
can cancel; condition (N) excludes this. The pole orders of the $v_i$ are
illustrative, not those of the computed section: for $a<M'$ any choice gives
distinct residues.}
\label{fig:residues}
